\chapter{Влияние оптической решётки на частоты переходов}\label{ch:ch2}

В конце прошлой главы мы пришли к выводу, что необходима разработка нового протокола начальных состояний часового перехода в оптической решётке, включающего охлаждение до основного колебательного состояния. 
Во введении рассматривались основные методы такой подготовки, и для этих целей возможно использовать переход на длине волны \thd.
Его спектральная ширина позволяет во-первых, селективно возбуждать переходы между различными магнитными подуровнями, а во-вторых, работать в режиме разрешённых колебательных частот. 
Но для этого требуется охарактеризовать влияние излучения оптической решётки на частоты как переходов между различными магнитными подуровнями основного и возбуждённого уровней перехода, так и между различными колебательными подуровнями выбранных магнитных подуровней. 
Данная глава посвящена исследованию влияния оптической решётки на длине волны $1063.5\unit{нм}$ на частоты компонент исследуемого перехода.

\section{Гамильтониан системы}
В этом пункте анализируется гамильтониан взаимодействия атомов с электрическим полем оптической решётки в присутствии внешнего задающего магнитного поля и описывается общий вид собственных состояний такой системы.
Магнитное поле $B$ предполагается постоянным и однородным, а электрическое поле оптической решётки записывается в виде $\mathbf{E} = \mathbf{E}^{(+)}\mathrm{exp}(-i \omega t) + \mathbf{E}^{(-)}\mathrm{exp}(i \omega t)$.
Общий вид гамильтониана взаимодействия атома с электрическим и магнитным полями такой конфигурации задаётся формулой \ref{eq:full_ham}.

\begin{equation} \label{eq:full_ham}
	\hat{H} = \hat{H}_{hf} + \left(\mu_B g_J \hat{J}_\mu + \mu_n g_I \hat{I}_\mu \right)B_\mu -\hat{\alpha}_{\mu\nu}E_{\mu}^{(-)}E_{\nu}^{(+)}\,,
\end{equation}

В этом уравнении $H_{hf}$ - гамильтониан, описывающий сверхтонкое расщепление, в случае атома тулия, ядерный спин которого $I=1/2$, представимый в виде $\hat{H}_{hf} = A_\textrm{hf}\,\hat{I}_\mu\,\hat{J}_\mu$.
Второе слагаемое описывает эффект Зеемана, где $g_J$ и $g_I$ - электронный и ядерный факторы Ланде.
Третье слагаемое описывает динамический эффект Штарка, параметром в котором является тензором поляризуемости $\alpha$.
Данный тензор раскладывается на тензоры 0, 1 и 2-го ранга, которые обычно называются скалярной, векторной и тензорной поляризуемостями по формуле \ref{eq:tensor}.
В этом уравнении $\alpha_s$, $\alpha_v$ и $\alpha_t$ -- скалярная, векторная и тензорная поляризуемости, $\delta_{\mu\nu}$ -- тензор Кронекера, а $\epsilon_{\xi\mu\nu}$ -- тензор Леви-Чевиты.
Суммирование ведётся по индексам $\xi \in \{x,y,z\}$. 

\begin{equation} \label{eq:tensor}
	\hat{\alpha}_{\mu\nu} = \alpha_s\delta_{\mu\nu}
	+\alpha_v i \epsilon_{\xi\mu\nu}\frac{\hat{J}_\xi}{J}
	+\alpha_t\frac{3}{J(2J-1)}\left[ \frac{1}{2}(\hat{J}_\mu \hat{J}_\nu + \hat{J}_\nu \hat{J}_\mu) - \frac{1}{3}\hat{\mathbf{J}}^2\delta_{\mu\nu} \right]\,,
\end{equation}

В случае присутствия только магнитного поля $B$ оператор полного момента $\hat{F} = \hat{J} + \hat{I}$ коммутирует с гамильтонианом, и собственными состояниями такой системы являются состояния с определённой проекцией полного момента $\ket{F, m_F}$.
В случае же учёта электрического поля решётки, оператор полного момента атома может коммутировать с гамильтонианом только в некоторых частных случаях.
В качестве примера такого частного случая служит часто встречающаяся в оптических часах конфигурация эксперимента: поляризация оптической решётки линейна и сонаправленна задающему магнитному полю.
В остальных же случаях возможно численно диагонализовать гамильтониан и разложить его собственные состояния (($\hat{H}\ket{\widetilde{m}_F}=E_{m_F}\ket{\widetilde{m}_F}$) \cite{virgo2013simultaneous}) в базисе состояний с определённой проекцией полного момента $\ket{\widetilde{m}_F} = \sum_{-F}^{F}c_{m_F} m_F$, $\ket{m_F}=\ket{F, m_F}$.
В данной работе эта задача решается с помощью пакета QuTip \cite{lambert2024qutip5quantumtoolbox} языка программирования Python.

\section{Подмешивание состояний}
В этом пункте рассматриваются собственные состояния гамильтониана \ref{eq:full_ham} в случае различных условий эксперимента для разных электронных уровней энергии.
Наиболее интересной для нас конфигурация эксперимента является случай линейно-поляризованной оптической решётки, электрическое поле которой перпендикулярно задающему магнитному полю.
Для такой конфигурации полей и электронных уровней, параметры которых нам известны -- основное и возбуждённое состояния часового перехода $\ket{g, F = 4(3)}$ и $\ket{c, F = 3, 2}$. %Насколько вообще можно записывать так уровень энергии в оптической решётке?
Поскольку рассматривается линейно-поляризованная оптическая решётка, то вклад векторной поляризуемости равен нулю. 
Значение скалярных и тензорных поляризуемостей, а также g-факторы рассматриваемых состояний приведены в таблице \ref{tab:StateMixingParameters}. 
Моделирование проводилось для магнитного поля $B = 340\unit{мГс}$ и глубины оптической решётки $U = 1000\Erec$. 
Результаты моделирования представлены на рисунке \ref{fig:StateMixing}.

\begin{table}[]
	\centering
	\begin{tabular}{|c|c|c|c|}
		\cline{1-4} 
		Уровень 			& $\alpha_s$ 	& $\alpha_t$	& g-фактор \\
		\gF{4} 				&  	1			& 1				&	1 		\\
		\gF{3} 				&  	1			& 1				&	1 		\\
		\clF{3} 			&  	1			& 1				&	1		\\
		\clF{2} 			&  	1			& 1				&	1 		\\
		\cline{1-4} 
	\end{tabular}
	\caption{Значения компонент динамической поляризуемости и g-факторов для основных и возбуждённых уровней часового перехода.}
	\label{tab:StateMixingParameters}
\end{table}

%Расчитываются собственные состояния основного и часового уровней в атомах тулия, а также возможный сдвиг частоты часового перехода, вызываемый этим эффектом.
Для характерных параметров нашего эксперимента, собственные состояния Гамильтониана практически идентичны состояниям с определённой проекцией магнитного момента, к которым примешивается незначительная доля состояний с другой проекцией.
Масштаб такого примешивания зависит от тензорной поляризуемости рассматриваемого уровня, а также от соотношения между величиной индукции задающего магнитного поля $B$ и электрического поля $E$. 
Для основного и часового уровней при характерных параметрах эксперимента доля подмешанных состояний не превышает долей процента, однако для уровней с большей тензорной поляризуемостью примесь состояний может достигать нескольких десятков процентов в зависимости от параметров эксперимента (см. рисунок \ref{fig:StateMixing}).

Построенная выше модель предсказывает различные сдвиги частоты микроволнового перехода между состояниями \gF{4} и \gF{3} для случая перпендикулярного и параллельного магнитных полей.
Такой сдвиг был обнаружен в работе \cite{mishin2024combined}.
Также были измерены зависимости величины такого сдвига от глубины оптической решётки для разных величин задающего магнитного поля (см. рисунок \ref{fig:StateMixingMW}). 

\begin{figure}[htbp!]
	\begin{center}
		\includegraphics[width=0.8\linewidth]
		{Chapter2/StateMixingMW}
	\end{center}
	\caption{ 
		Зависимость сдвига частоты микроволнового перехода \MW от глубины оптической решётки для различных магнитных полей. 
		Точками показаны экспериментальные данные (?), пунктирной линией модель, залита 1$\sigma$.
	}
	\label{fig:StateMixingMW}
\end{figure}

 
Для синтетической частоты часового перехода этот эффект тоже может быть наблюдаем: предсказание модели на рисунке \ref{fig:StateMixingSynthetic}.
Проверка этого расчёта затруднена, т.к. неотделима от измерения поляризуемостей, однако предсказание сдвига до $1\unit{Гц}$ впечатляет.

\begin{figure}[htbp!]
	\begin{center}
		\includegraphics[width=0.8\linewidth]
		{Chapter2/StateMixingSynthetic}
	\end{center}
	\caption{ 
		Сдвиг синтетической частоты часового перехода из-за подмешивания состояний в зависимости от глубины оптической решётки и величины магнитного поля в логарифмическом масштабе.
	}
	\label{fig:StateMixingSynthetic}
\end{figure}

Для работы с атомами тулия в оптических пинцетах также интересен случай дипольной ловушки на длине волны $532\unit{нм}$ с различными параметрами эллиптичности формирующего дипольную ловушку излучения.
Для основного состояния \gF{4}, векторная поляризуемость на длине волны $532\unit{нм}$ равна $\alpha_v = 1$, остальные параметры указаны в таблице \ref{tab:StateMixingParameters}. 
Результаты моделирования для линейной (перпендикулярной магнитному полю $B$), немного-эллиптичной, а также круговой поляризации показаны на рисунке \ref{fig:StateMixingElliptical}.
Виден абсолютный бардак в разложении собственных состояний такой системы по магнитным подуровням.

\begin{figure}[htbp!]
	\begin{center}
		\includegraphics[width=0.8\linewidth]
		{Chapter2/StateMixingElliptical.pdf}
	\end{center}
	\caption{ 
		Разложение собственных состояний гамильтониана \ref{eq:full_ham} по состояниям с определённой проекцией магнитного момента для уровня \gF{4} и поляризаций дипольной ловушки разной степени эллиптичности.
	}
	\label{fig:StateMixingElliptical}
\end{figure}

\section{Поляризуемость охлаждающего перехода}

Данный раздел посвящен характеризации влияния оптической решётки на частоту охлаждающего перехода.
В формуле \ref{eq:full_ham}, параметрами являются g-факторы, а также поляризуемость исследуемого уровня. 
Факторы Ланде основного и возбуждённого состояний охлаждающего перехода были известны \cite{}, как и компоненты динамической поляризуемости основного состояния на длине волны \lt \cite{}, а значит, оставалось измерить компоненты динамической поляризуемости возбуждённого состояния.

Для измерения поляризуемости исследовалась зависимость частот переходов между различными Зеемановскими компонентами основного и возбуждённого состояний охлаждающего перехода. 
Ранее были разработаны протоколы подготовки атомов тулия на центральном и крайнем магнитном подуровне основного состояния \cite{}, соответственно для исследования было доступно 6 различных переходов между Зеемановскими подуровнями. 
Возбуждение исследуемой компоненты охлаждающего перехода выполнялось пучками МОЛ коротким немощным (параметр насыщения $s\approx1$) импульсом.
Для измерения итоговой населённости магнитных подуровней выполнялся эксперимент Штерна-Герлаха: после подготовки атомов на начальном магнитном подуровне и возбуждении исследуемого перехода включалось магнитное поле с большим градиентом вдоль оси Z, удерживающее поле оптической решётки выключалось, и через $20\unit{мс}$ делалась фотография облака, разделившегося на несколько атомных ансамблей с разной проекцией магнитного момента на ось квантования.
Полученное изображение аппроксимировалась суммой Гауссиан, и ``амплитуда'' соответствующей Гауccианы выбиралась в качестве населённости соответствующего подуровня.
%Так как эффект подмешивания состояний для основного состояния мал (его масштаб составляет мене 1\%), его влиянием в данном эксперименте пренебрегаем. 


Измеренная зависимость доли деполяризованных атомов от частоты возбуждающего излучения аппроксимировалась функцией Гаусса, центр которой принимался за центральную частоту перехода при исследуемых параметрах. 
Переход между подуровнями $\ket{g, m_F = -4} \xrightarrow{} \ket{e, m_F-5}$ не вызывал деполяризации атомов, поэтому определить центральную частоту этого перехода было возможно только по потерям атомов из оптической решётки в результате нагрева, что сопряжено с большими ошибками измерения, поэтому исследовались частоты пяти оставшихся переходов в зависимости от глубины решётки. 
Результат аппроксимировался набором прямых, наклоны которых аппроксимировались построенной выше моделью, параметрами которой являлись скалярная и тензорная поляризуемость уровня \eF{5}.
Результаты представлены на рисунке \ref{fig:Depolarization}.
Полученные таким образом значения скалярной и тензорной поляризуемостей оказались равны $\alpha_s = 190(30)\unit{ат.ед.}$ и $\alpha_t = 47(7)\unit{ат.ед.}$.
Погрешности абсолютных значений связаны с неопределённостью поляризуемости основного состояния, для определения которого необходимо точно измерять интенсивность излучения, формирующего оптическую решётку.
Для исследования спектра магнитных подуровней перехода $506.2\unit{нм}$ в этой работе нас интересуют только относительные величины $\bar{\alpha}$.

\begin{figure}[htbp!]
	\begin{center}
		\includegraphics[width=0.8\linewidth]
		{Chapter2/Depolarization.pdf}
	\end{center}
	\caption{
		Описание эксперимента по измерению компонент динамических поляризуемостей на длине волны $1063.5\unit{нм}$ верхнего уровня охлаждающего перехода с длиной волны $506.2\unit{нм}$.
		а) Эксперимент Штерна-Герлаха и населённости магнитных после деполяризации чистого состояния
		б) Пример спектра перехода, полученного с помощью деполяризации чистого состояния
		в) Набор прямых, полученных в результате аппроксимации экспериментальных данных теорией.
		Чёрным цветом показана полученная из модели частота перехода $\ket{g, F = -4} \xrightarrow{} \ket{e, F = -5}$.
		Видно, что она практически не смещается оптической решёткой. 
	}
	\label{fig:Depolarization}
\end{figure}

Из измеренных значений компонент поляризуемостей можно сделать несколько выводов.
Во-первых, частота охлаждающего перехода $\ket{g, m_F = -4} \xrightarrow{} \ket{e, m_F = -5}$ практически не смещается излучением решётки (см. рис. \ref{fig:Depolarization}): для типичного значения глубины $ U = 1000 E_{rec}$ и магнитного поля $B = 600 \unit{мГс}$ сдвиг частоты составлял $\delta \nu = 6\unit{кГц}$, что меньше естественной ширины линии охлаждающего перехода. 
Это позволяет эффективно охлаждать атомы тулия на этом переходе (подробнее см. в следующей главе). 

Во-вторых, для возбуждённого состояния охлаждающего перехода эффект подмешивания состояний оказывается достаточно большим: вклад неосновных магнитных подуровней может составлять до 1\% при характерных параметрах эксперимента.
Это, и то, что спектральная ширина перехода меньше Зеемановского расщепления при используемых в эксперименте магнитных полях, позволяет зарегистрировать компоненты перехода, запрещённые в дипольном приближении (см. рис. \ref{fig:StateMixingTransition}).
Так то, что к состоянию \emF{5}{3} примешивается состояние \emF{5}{1} делает отличным от нуля матричный элемент $\bra{g, F = 4, m_F = 0}\ket{e, F = 5, m_F = 3}$, что приводит к пику в спектре перехода на соответствующей частоте и характерному для распада с состояния \emF{5}{3} распределению атомов после его возбуждения.

Измеренные значения поляризуемостей позволяют выделить два различных режима работы с использованием перехода $506.2\unit{нм}$.
При  большом значении магнитного поля ($B > 500\unit{мГс}$) частоты переходов определяются в основном эффектом Зеемана, а сдвиг из-за эффекта Штарка является лишь малой поправкой, соответственно, картина в целом не меняется. 
При использовании меньших значений магнитного поля ($B < 200\unit{мГс}$), вклад эффекта Штарка в сдвиг частоты переходов становится соизмерим с вкладом эффекта Зеемана.
Промежуточные значения магнитного поля в эксперименте не использовались из-за наличия резонанса Фешбаха для состояния $\ket{g, F = 4, m_F = -4}$ в магнитном поле $B \approx 430\unit{мГс}$ \cite{khlebnikov2019random}, приближение к которому  вызывает значительные потери атомов в процессе охлаждения.  

\begin{figure}[htbp!]
	\begin{center}
		\includegraphics[width=0.8\linewidth]
		{Chapter2/StateMixingClock.pdf}
	\end{center}
	\caption{ 
		Разложение собственных состояний гамильтониана \ref{eq:full_ham} по состояниям с определённой проекцией магнитного момента. 
		Расчет представлен для состояний $\ket{g, F = 4}$, $\ket{c, F = 3}$, и $\ket{e, F = 5}$.
		Чёрным цветом выделены основное и возбуждённое состояние часового перехода. 
		Видно, что к состоянию с нулевой проекцией полного момента примешиваются также состояния с другими проекциями, что вызывает дополнительную чувствительность часового перехода к магнитному полю.
		Рассчёт для уровня $\ket{e, F = 5}$ выполнялся на основе измеренных в этой главе компонент динамической поляризуемостей этого состояния.
	}
	\label{fig:StateMixing}
\end{figure} 
 
\begin{figure}[htbp!]
	\begin{center}
		\includegraphics[width=0.8\linewidth]
		{Chapter2/StateMixingTransition.pdf}
	\end{center}
	\caption{
		Observation of $|\Delta m_F| > 1$ lines one electric-dipole transition at 506\,nm in neutral thulium atoms in optical lattice.
		Configurations of the experiments (a)-(c). 
		Measured excitation spectra of $\ket{g,m_F=0} \rightarrow \ket{e, m_F}$ transitions (d).
		Horizontal axis is in the units $g_e\mu_B B_0\sim500$\,kHz of Zeeman splitting of the upper level, $0$ corresponds to $\ket{g,m_F=0} \rightarrow \ket{e, m_F=0}$.
		Orange dots represent data measured in configuration (a) when the lattice polarization $\vec{\varepsilon}_L$ is perpendicular to the bias magnetic field $\vec{B}_0$ and excitation radiation has $\sigma^\pm$ components.
		Besides expected $\ket{g,m_F=0} \rightarrow \ket{e, m_F=\pm 1}$ transitions one can clearly see peak at $\ket{g,m_F=0} \rightarrow \ket{e, m_F=3}$ transition detuning.
		Blue dots illustrate spectrum measured in configuration (b) when (similar to the previous case $\vec{\varepsilon}_L \perp \vec{B}_0$ but radiation has only $\pi$ polarization.
		Yet again, besides $\ket{g,m_F=0} \rightarrow \ket{e, m_F=0}$ peak we observe excitation a detuning corresponding to $\ket{g,m_F=0} \rightarrow \ket{e, m_F=2}$ transition.
		Blue points show excitation spectrum measured when the lattice polarization and bias magnetic field are parallel.
		No unexpected features were observed in this configuration.
		Level diagram (e) shows excitation of $\ket{g,m_F=0} \rightarrow \ket{e, m_F=3}$ transition and expected decay channels from $\ket{e, m_F=3}$.
		Measured population distribution if $\ket{g}$ manifold is show in orange dots in (f), grey bars indicate probabilities of spontaneous decay from $\ket{e, m_F=3}$. 
		Plots (g) and (h) show measured amplitude of the $\ket{g,m_F=0} \rightarrow \ket{e, m_F=2}$ ($\ket{g,m_F=0} \rightarrow \ket{e, m_F=3}$) peaks as a function of lattice depth. 
		Dashed lines are parabolic fits.
	}
	\label{fig:StateMixingTransition}
\end{figure}

\section{Основные результаты главы \ref{ch:ch2}}.

\begin{itemize}
	\item Обнаружен и численно охарактеризован эффект подмешивания состояний в случае непараллельных электрического поля оптической решётки и задающего ось квантования магнитного поля, для атомов тулия, удерживающихся в потенциале оптической решётки.
	Показано, что этот эффект может приводить к возбуждению компонент перехода, запрещённых в дипольном приближении, а также сдвигу частоты часового перехода величиной до $1\unit{Гц}$ для магнитного поля $B = 100\unit{мГс}$ и глубины оптической решётки $U = 1000 \Erec$.
	\item Измерены динамические скалярная и тензорная поляризуемости  на длине волны $1063.5\unit{нм}$ возбуждённого состояния перехода на длине волны $506.2\unit{нм}$, которые равны соответственно $\alpha_s = 190(30)~\text{ат.ед}$ и $\alpha_t = 47(7)~\text{ат.ед}$.
\end{itemize}